\subsection{整闭代数中的纯量扩张}

命题 19. 设 \(k\) 是域，\(L\) 是 \(k\) 的可分扩张，\(R\) 是整闭的 \(k\) -代数。若环 \(L \otimes_k R\) 是整环，则它是整闭的。

设 K 是 R 的分式域；由于 k 是域，\(L \otimes_{k} R\) 典范地等同于 \(L \otimes_{k} K\) 的子 k-代数，而 L 与 R 典范地等同于 \(L \otimes_{k} R\) 的子 k-代数。此外，由于 R 的元素 \(s \neq 0\) 不是 R 中的零因子，且 L 在 k 上平坦（第 I 章，§2，第 3 节），故 \(1 \otimes s\) 不是 \(L \otimes_{k} R\) 中的零因子；把 s 与 \(1 \otimes s\) 等同，于是可以看出，若 S = R - \{0\}，则 \(L \otimes_{k} K\) 等同于 \(S^{-1}(L \otimes_{k} K)\) ；由于假设 \(L \otimes_{k} R\) 是整环，\(L \otimes_{k} K\) 便等同于 \(\Omega\) ——\(L \otimes_{k} R\) 的分式域——的子环。

\begin{enumerate}
\def\labelenumi{(\arabic{enumi})}
\tightlist
\item
  假设 \(\mathbf{L}\) 是 \(k\) 的有限扩张；则 \(\mathbf{L} \otimes_k \mathbf{K}\) 是 \(\mathbf{K}\) 上的有限秩代数，且由假设没有零因子；因此它是域（《代数》第 V 章，§2，第 1 节，命题 1），从而在此情形它就是 \(\Omega\) ，即 \(\mathbf{L} \otimes_k \mathbf{R}\) 的分式域。设 \((w_1, \ldots, w_n)\) 是 \(\mathbf{L}\) 在 \(k\) 上的基，它因而也是 \(\mathbf{L} \otimes_k \mathbf{K}\) 在 \(\mathbf{K}\) 上的基。存在一组基 \((w_1^*, \ldots, w_n^*)\) （\(\mathbf{L}\) 的基），使得 \(\mathrm{Tr}_{\mathrm{L/k}}(w_i w_j^*) = \delta_{ij}\) （第 6 节，引理 3）；每个 \(z \in \mathbf{L} \otimes_k \mathbf{K}\) 都可唯一地写成 \(z = \sum_{i=1}^{n} a_i w_i\) ，其中 \(a_i \in \mathbf{K}\) ；于是
\end{enumerate}

\[
\operatorname{Tr} _ {(\mathrm{L} \otimes \mathrm{K}) / \mathrm{K}} (z w _ {j} ^ {*}) = \sum_ {i = 1} ^ {n} a _ {i} \operatorname{Tr} _ {(\mathrm{L} \otimes \mathrm{K}) / \mathrm{K}} (w _ {i} w _ {j} ^ {*})
\]

又由于在 L 中迹 \(\mathrm{Tr}_{(L\otimes K)/K}\) 与 \(\mathrm{Tr}_{L/K}\) 一致（《代数》第 VIII 章，§ 12，第 2 节，公式 (13)），最终得 \(\mathrm{Tr}_{(L\otimes K)/K}(zw_j^*) = a_j\) ，\(1 \leqslant j \leqslant n\) 。另一方面注意 L 的元素在 \(k\) 上是整的，从而也在 R 上是整的（第 1 节，命题 2 的推论 1）；因此（第 1 节，命题 5）\(\mathbf{L} \otimes_k \mathbf{R}\) 在 R 上是整的。在这样的约定下，设 \(z \in \mathbf{L} \otimes_k \mathbf{K}\) 在 \(\mathbf{L} \otimes_k \mathbf{R}\) 上是整的；则 \(z\) 也在 R 上是整的（第 1 节，命题 6），从而 \(zw_j^*\) 亦然，于是 \(a_j = \mathrm{Tr}_{(L\otimes K)/K}(zw_j^*)\) （\(1 \leqslant j \leqslant n\)）也在 R 上是整的（第 6 节，命题 17 的推论 2）。由于 R 是整闭的，\(a_j \in \mathbf{R}\) 对所有 \(j\) 成立，从而 \(z \in \mathbf{L} \otimes_k \mathbf{R}\) ，这就证明了此情形下的命题。

\begin{enumerate}
\def\labelenumi{(\arabic{enumi})}
\setcounter{enumi}{1}
\tightlist
\item
  现在假设 \(\mathbf{L}\) 是 \(k\) 的有限生成的可分扩张；则存在 \((x_1, \ldots, x_d)\) ，它是 \(\mathbf{L}\) 在 \(k\) 上的分离超越基（《代数》第 V 章，§ 9，第 3 节，定理 2）；由于 \(\mathbf{L}\) 与 \(\mathbf{K}\) 在 \(k\) 上于域 \(\Omega\) 中代数不相交（《代数》第 V 章，§ 5，第 4 节），故诸 \(x_i\) 在 \(\mathbf{K}\) 上代数无关；从而 \(\mathrm{R}[x_1, \ldots, x_d]\) 是整闭的（第 3 节，命题 13 的推论 2）。设 \(\mathbf{T}\) 是由 \(\neq 0\) 的元素组成的集（环 \(\mathbf{A} = k[x_1, \ldots, x_d] \subset \mathbf{L}\) 的），于是域 \(k_1 = k(x_1, \ldots, x_d) \subset \mathbf{L}\) 等于 \(\mathbf{T}^{-1} k[x_1, \ldots, x_d]\) ；则
\end{enumerate}

\[
\begin{array}{r l} k _ {1} \otimes_ {k} \mathrm{R} & = (\mathrm{T} ^ {- 1} \mathrm{A}) \otimes_ {k} \mathrm{R} = \mathrm{T} ^ {- 1} \mathrm{A} \otimes_ {\mathrm{A}} (\mathrm{A} \otimes_ {k} \mathrm{R}) = \mathrm{T} ^ {- 1} (\mathrm{A} \otimes_ {k} \mathrm{R}) \\ & = \mathrm{T} ^ {- 1} \mathrm{R} [ x _ {1}, \dots , x _ {d} ] \end{array}
\]

（由张量积的结合律），因此这个整环是整闭的（第 5 节，命题 16 的推论 1）。但 \(\mathbf{L} \otimes_k \mathbf{R}\) 等同于 \(\mathbf{L} \otimes_{k_1} (k_1 \otimes_k \mathbf{R})\) ，且按定义 \(\mathbf{L}\) 是 \(k_1\) 的有限可分扩张；于是由 (1) 可得 \(\mathbf{L} \otimes_k \mathbf{R}\) 是整闭的。

\begin{enumerate}
\def\labelenumi{(\arabic{enumi})}
\setcounter{enumi}{2}
\tightlist
\item
  一般情形。若 z 是 \(\Omega\) 中在 \(L \otimes_{k} R\) 上整的元素，则它满足形如 \(z^{m} + b_{1}z^{m-1} + \cdots + b_{m} = 0\) 的关系式，其中诸 \(b_{i}\) 属于 \(L \otimes_{k} R\) ；于是存在 L 在 k 上的有限生成子扩张 \(L'\) ，使得诸 \(b_{i}\) 属于 \(L' \otimes_{k} R\) （\(1 \leqslant i \leqslant m\)），且 z 属于 \(L' \otimes_{k} K\) 。于是由 (2) 得 \(z \in L' \otimes_{k} R\) ，从而 \(L \otimes_{k} R\) 是整闭的。
\end{enumerate}

* 设 V 是不可约仿射代数簇，k 是 V 的定义域，R 是 k 上定义的 V 上的正则函数所成的环；若 R 是整闭的，则称 V 在 k 上是正规的；命题 19 表明，若 V 在 k 上正规，则它在 k 的每个可分扩张 L 上仍是正规的。*

推论. 设 \(k\) 是域，\(\mathbf{R}\) 与 \(\mathbf{S}\) 是两个整闭的 \(k\) -代数。假设环 \(\mathbf{R} \otimes_k \mathbf{S}\) 是整环，且分式域 \(\mathbf{K}\) 与 \(\mathbf{L}\) （分别为 \(\mathbf{R}\) 与 \(\mathbf{S}\) 的分式域）在 \(k\) 上可分。则环 \(\mathbf{R} \otimes_k \mathbf{S}\) 是整闭整环。

由于 R 与 S 等同于 \(R \otimes_{k} S\) 的子代数，K 与 L 等同于 \(\Omega\) ——\(R \otimes_{k} S\) 的分式域——的在 k 上线性不相交的子域（《代数》第 V 章，§2，第 3 节，命题 5）。于是由命题 19 可知 \(R \otimes_{k} L\) 与 \(K \otimes_{k} S\) 都是整闭整环；由于它们的交是 \(R \otimes_{k} S\) （第 I 章，§2，第 6 节，命题 7），故 \(R \otimes_{k} S\) 是整闭整环（第 2 节，命题 8 的推论）。

* 给定 k 上定义的两个不可约仿射簇 V，W，它们的乘积 \(V \times W\) 是仿射簇，且 \(V \times W\) 上的正则函数环等同于 V 上正则函数环与 W 上正则函数环在 k 上的张量积。命题 19 的推论表明，若 V 与 W 都在 k 上正规，则 \(V \times W\) 在 k 上正规。*

